\chapter{기술통계 --- 숫자로 데이터를 요약하기}
\label{ch:descriptive}

\begin{inbrief}
  \item[Why] 데이터를 몇 개의 숫자로 줄이려면, 먼저 ``값이 나올 가능성''을
             적는 언어인 확률변수(Random Variable), 확률질량함수(PMF),
             확률밀도함수(PDF), 누적분포함수(CDF)가 필요하다.
  \item[When] 새 데이터를 받아 첫 요약표를 만들 때, 그리고 분포 모형이나
              가설검정(Hypothesis Testing)으로 넘어가기 전에 쓴다.
  \item[Where] 주사위, 붓꽃(Iris) 꽃받침 길이(Sepal Length), 과제 데이터의
               \texttt{sugars}(당류) 변수에 적용한다.
  \item[How] 확률변수로 값을 정의하고, PMF/PDF와 CDF로 분포를 적은 뒤,
             중심(Center, 평균, 중앙값, 최빈값(Mode))과
             퍼짐(Spread, 범위(Range), 표준편차, IQR)을
             저항성(Resistance) 기준으로 고른다.
\end{inbrief}
% source: deck 4 slide 1

\term{기술통계}(Descriptive Statistics)는 데이터를 몇 개의 숫자로 줄이는
도구다. \analogy{지도}를 그리는 일이다(\Cref{tab:analogy}). 줄인다는 것은
반드시 무언가를 버린다는 뜻이고, 무엇을 버렸는지 모르면 남은 숫자를 잘못 읽는다.

이 장은 네 가지 요약값을 다룬다. 중심(Center), 퍼짐(Spread), 비대칭(Skewness),
그리고 값이 아예 없는 경우(Missing Value)다.

\begin{keyidea}
요약통계와 그림은 \textbf{한 쌍으로 다녀야 한다}. 요약값은 ``무엇이 일어났는가''
를 알려 주고, 그림은 ``왜 그런 값이 나왔는가''를 알려 준다. 둘 중 하나만으로
내린 결론은 이 장에서 계속 반복해서 틀린다.
\end{keyidea}

% =====================================================================
\section{배경}
\label{sec:desc-background}
% source: deck 4 slide 1-3

이 장의 출처는 \texttt{notes/4\_BasicStats.pptx}(23 슬라이드) 하나다.
슬라이드가 외부 문헌을 밝히지 않는 자리는 슬라이드 자체를 출처로 적는다.

기술통계를 계산하려면 그 전에 데이터를 어떻게 표로 담고, 각 값이 얼마나
자주 나오는지 어떻게 적을지부터 정해야 한다. 이 장은 다음 두 가지를
그대로 가져와 쓴다.

\begin{itemize}
  \item 속성(Attribute)과 척도: 값이 숫자인지 범주인지, 순서가 있는지를
        가른다(\Cref{ch:attributes}, \Cref{sec:attr-numeric}).
  \item 데이터 행렬(Data Matrix) $n \times d$: 행이 객체, 열이 속성인 표
        (\Cref{sec:data-matrix}).
\end{itemize}

이 장에서 계속 쓰는 데이터는 붓꽃(Iris) 데이터셋이다. 150송이의 붓꽃마다
꽃받침 길이(Sepal Length)를 포함한 네 가지 속성을 측정했고, 꽃받침
길이는 4.3\,cm에서 7.9\,cm 사이에 분포한다.

여러 속성 중 하나씩만 골라 한 번에 한 속성만 들여다보는 일을
\term{일변량 분석}(Univariate Analysis)이라고 부른다.
% source: deck 4 slide 14
이때 표의 한 열, 즉 속성 하나가 이 장에서 다룰 확률변수 하나에
대응한다.

이 장의 비유는 \Cref{ch:descriptive}의 \analogy{지도}를 그대로 잇는다.
확률분포(Probability Distribution)는 지형 전체이고, 이 장에서 배우는
평균, 중앙값, 표준편차 같은 요약값은 그 지형 위에 찍은 몇 개의 표지판이다.
표지판만 보고 지형 전체를 상상할 수 없듯, 요약값 몇 개만으로 분포
전체를 복원할 수는 없다.

% =====================================================================
\section{확률변수}
\label{sec:random-variable}
% source: deck 4 slide 4-5

\begin{inbrief}
  \item[Why] 관측값을 계산에 넣을 수 있는 숫자로 바꾸는 규칙이 필요하다.
  \item[When] 속성을 확률 계산이나 수식에 넣기 전에 정의한다.
  \item[Where] 수치형 속성은 있는 그대로 쓰고, ``조건을 만족하는가''
               같은 판정은 0과 1로 바꾼다.
  \item[How] 표본공간의 각 결과에 실수 하나를 대응시킨다.
\end{inbrief}

\paragraph{고치는 문제.} ``이 붓꽃은 꽃받침이 길다''라는 말은 그 자체로는
계산할 수 없다. 평균을 내거나 확률을 구하려면 먼저 이 말을 숫자로 바꿔야
한다.

\begin{definition}[확률변수]
\label{def:random-variable}
\term{표본공간}(Sample Space) $\mathcal{O}$는 실험에서 나올 수 있는 모든
결과의 집합이다. \term{확률변수}(Random Variable) $X$는 표본공간의 각
결과를 실수 하나로 대응시키는 함수다.
\[
  X : \mathcal{O} \to \mathbb{R}
\]
$X$가 유한 개 또는 셀 수 있는(Countably Infinite) 값만 가지면
\term{이산 확률변수}(Discrete Random Variable), $\mathbb{R}$의 어떤
구간에서도 값을 가질 수 있으면 \term{연속 확률변수}(Continuous Random
Variable)라고 부른다.
\end{definition}

수치형 속성 $X_j$는 별다른 변환 없이도 이미 확률변수다. 결과 $v$를
그 값 자신으로 보내는 \term{항등 확률변수}(Identity Random Variable)
\[
  X(v) = v, \qquad v \in \mathcal{O} = \mathbb{R}
\]
이기 때문이다.

\begin{example}[꽃받침 길이를 0/1로 바꾸기]
\label{ex:iris-long-sepal}
``꽃받침이 길다''를 ``꽃받침 길이가 7\,cm 이상이다''로 정하면, 확률변수
$A$를
\[
  A(v) = \begin{cases} 1 & v \ge 7 \\ 0 & v < 7 \end{cases}
\]
로 정의할 수 있다. $A$가 가지는 값은 $\{0,1\}$뿐이므로 $A$는 이산
확률변수다.
\end{example}

연속 확률변수의 예(삼각형 모양 밀도)는 \Cref{sec:pmf-pdf}의
\Cref{ex:triangle-pdf}에서 이어 다룬다. 한 절에서는 한 개념만 새로
들여온다.

표본공간을 적을 때 주의할 점이 있다. \Cref{def:random-variable}은
수치형 속성의 표본공간을 $\mathcal{O} = \mathbb{R}$이라고 적지만, 실제로
관측되는 꽃받침 길이는 $[4.3, 7.9]$ 구간에만 있다. $[4.3, 7.9]$는
\textbf{이번에 관측된 범위}이고, 표본공간은 정의상 \textbf{나올 수 있는
모든 결과}다. 관측 범위가 좁다고 표본공간 자체가 좁아지는 것은 아니다.

% =====================================================================
\section{확률질량함수와 확률밀도함수}
\label{sec:pmf-pdf}
% source: deck 4 slide 6-8, 15

\begin{inbrief}
  \item[Why] 확률변수의 각 값이 얼마나 자주 나오는지를 적는 함수가
             필요하다.
  \item[When] 이산 확률변수는 PMF, 연속 확률변수는 PDF를 쓴다.
  \item[Where] 주사위 눈, 붓꽃의 0/1 판정에는 PMF를, 꽃받침 길이 같은
               연속 측정값에는 PDF를 쓴다.
  \item[How] PMF는 값마다의 확률을 더하고, PDF는 구간에서 적분한다.
\end{inbrief}

\paragraph{고치는 문제.} 연속 확률변수에서는 한 점 $X=v$가 나올 확률이
정확히 0이다(\Cref{def:pdf} 뒤에서 다시 본다). 그러므로 ``이 값이 나올
확률이 얼마인가''라는 질문이 이산형에서는 성립해도 연속형에서는 성립하지
않는다. 그래서 이산형과 연속형은 서로 다른 함수로 분포를 적는다.

\begin{definition}[확률질량함수]
\label{def:pmf}
이산 확률변수 $X$의 \term{확률질량함수}(Probability Mass Function, PMF)
$f$는
\[
  f(x) = P(X = x)
\]
로 정의되고, 다음 두 성질을 만족한다. $f(x) \ge 0$이고
$\sum_x f(x) = 1$이다.
\end{definition}

\Cref{ex:iris-long-sepal}의 $A$로 PMF를 계산해 보자. 150송이 중
꽃받침 길이가 7\,cm 이상인 것은 13송이다.

\begin{align*}
  f(1) &= P(A=1) = 13/150 = 0.0867  && \text{길다고 판정된 비율}\\
  f(0) &= P(A=0) = 137/150 = 0.9133 && \text{나머지 비율}
\end{align*}

$f(0) + f(1) = 1$로 \Cref{def:pmf}의 성질을 만족한다. $A$처럼 값이
$\{0,1\}$뿐인 확률변수의 분포를 \term{베르누이 분포}(Bernoulli
Distribution)라고 부른다. 이름 있는 분포로서의 자세한 성질은
\Cref{sec:binomial}에서 다시 다룬다.

\begin{definition}[확률밀도함수]
\label{def:pdf}
연속 확률변수 $X$의 \term{확률밀도함수}(Probability Density Function,
PDF) $f$는 구간 확률을 다음과 같이 넓이로 정의한다.
\[
  P(X \in [a,b]) = \int_a^b f(x)\,dx, \qquad \int_{-\infty}^{\infty} f(x)\,dx = 1
\]
한 점의 확률은 폭이 0인 구간의 넓이이므로 항상 $P(X=v) = 0$이다.
\end{definition}

\begin{example}[삼각형 모양의 밀도]
\label{ex:triangle-pdf}
\[
  f(x) = \begin{cases}
    x     & 0 \le x \le 1 \\
    2 - x & 1 < x \le 2 \\
    0     & \text{그 밖}
  \end{cases}
\]
그래프로 그리면 $(0,0)$-$(1,1)$-$(2,0)$을 잇는 삼각형이다
(\Cref{fig:triangle-pdf-cdf}). 넓이가 1인지 먼저 확인한다.

\begin{align*}
  &\tfrac{1}{2}\times 1 \times 1 = 0.5 && \text{왼쪽 삼각형(밑변 1, 높이 1) 넓이}\\
  &\tfrac{1}{2}\times 1 \times 1 = 0.5 && \text{오른쪽 삼각형 넓이}\\
  &0.5 + 0.5 = 1                       && \text{전체 넓이, \Cref{def:pdf}의 조건과 일치}
\end{align*}

구간 확률도 넓이로 구한다. $P(0.5 \le X \le 1.5)$는 전체 넓이 1에서
양끝의 작은 삼각형 두 개를 뺀 값이다. 양끝 삼각형은 각각 밑변 0.5,
높이 0.5다.

\begin{align*}
  &2 \times \left(\tfrac{1}{2}\times 0.5 \times 0.5\right) = 0.25 && \text{양끝 두 삼각형의 넓이 합}\\
  &1 - 0.25 = 0.75                                                 && \text{$P(0.5 \le X \le 1.5)$}
\end{align*}
\end{example}

\begin{figure}[ht]\centering
  \includegraphics[width=\linewidth]{04_02_triangle_pdf_cdf.png}
  \caption{삼각형 밀도(PDF, 왼쪽)와 그 누적분포함수(CDF, 오른쪽). 색칠한
  구간 $[0.5,1.5]$의 넓이가 0.75이고, CDF에서는 $F(1)=0.5$인 지점이
  표시되어 있다.}
  \label{fig:triangle-pdf-cdf}
\end{figure}

\begin{pitfall}[PDF 값 자체는 확률이 아니다]
\label{pit:pdf-not-prob}
$f(x)$가 1을 넘어도 PDF의 정의를 어기지 않는다. 예를 들어 폭 0.5인
구간에 고르게 퍼진 밀도는 $f(x) = 2$다($0.5 \times 2 = 1$로 넓이가
1이어야 하기 때문이다). 확률은 언제나 \Cref{def:pdf}처럼 \textbf{넓이}로
읽는다. 높이 하나만 보고 확률이라고 읽으면 틀린다.
\end{pitfall}

\subsection{기댓값}
\label{subsec:expectation}

\term{기댓값}(Expected Value) $\E[X]$는 확률변수가 평균적으로 어떤 값을
갖는지를 나타낸다. 이산형은 PMF로 가중 평균을 내고, 연속형은 PDF로
적분한다.
\[
  \E[X] = \sum_x x f(x) \qquad \text{(이산)}, \qquad
  \E[X] = \int_{-\infty}^{\infty} x f(x)\,dx \qquad \text{(연속)}
\]

주사위에서는 여섯 눈이 각각 $1/6$의 확률을 가지므로

\begin{align*}
  \E[X] &= \sum_{k=1}^{6} k \cdot \tfrac{1}{6} = \tfrac{1+2+3+4+5+6}{6}
    && \text{PMF로 가중 평균}\\
  \E[X] &= \tfrac{21}{6} = 3.5 && \text{약분}
\end{align*}

\Cref{ex:triangle-pdf}의 삼각형 밀도에서는

\begin{align*}
  \E[X] &= \int_0^1 x \cdot x\,dx + \int_1^2 x(2-x)\,dx
    && \text{구간별로 나눠 적분}\\
  \E[X] &= \tfrac{1}{3} + \tfrac{2}{3} && \text{각 적분값}\\
  \E[X] &= 1 && \text{합}
\end{align*}

두 계산 모두 도형의 무게중심과 일치한다. 표본에서 이 기댓값을
추정한 값이 \Cref{def:center}의 평균 $\hat\mu$다.

% =====================================================================
\section{누적분포함수}
\label{sec:cdf-def}
% source: deck 4 slide 9-11, 18

\begin{inbrief}
  \item[Why] 이산과 연속을 한 함수로 다루고, ``이 값 이하가 나올
             확률''을 바로 읽기 위해서다.
  \item[When] 분포를 비교하거나, 분위수(Quantile)나 중앙값을 구할 때
              쓴다.
  \item[Where] 주사위, 붓꽃 데이터의 이항·정규 근사, 뒤에서 다룰
               분포 비교와 가설검정에 쓰인다.
  \item[How] PMF는 누적합으로, PDF는 누적적분으로 CDF를 만든다.
\end{inbrief}

\paragraph{고치는 문제.} PMF와 PDF는 이산형과 연속형에서 서로 다른
종류의 함수라서 같은 축 위에 겹쳐 놓고 비교할 수 없다. 이산형과
연속형을 같은 잣대로 비교하려면, 둘 다 $[0,1]$ 사이의 값을 내놓는
함수가 필요하다.

\begin{definition}[누적분포함수]
\label{def:cdf}
확률변수 $X$의 \term{누적분포함수}(Cumulative Distribution Function,
CDF) $F$는
\[
  F(x) = P(X \le x), \qquad F : \mathbb{R} \to [0,1]
\]
로 정의된다. 이산형에서는 $F(x) = \sum_{u \le x} f(u)$이고, 연속형에서는
$F(x) = \int_{-\infty}^{x} f(u)\,du$다. $F$는 네 가지 성질을 항상
만족한다. 감소하지 않는다, 오른쪽에서 연속이다(Right-continuous, 즉
$\lim_{h \to 0^+} F(x+h) = F(x)$), $F(x) \to 0$ ($x \to -\infty$),
$F(x) \to 1$ ($x \to \infty$).
\end{definition}

\begin{example}[주사위 눈의 CDF]
\label{ex:dice-cdf}
\Cref{tab:dice-cdf}은 주사위 눈 $k$마다 $P(X=k)$와 $F(k)=P(X\le k)$를
정리한다(\Cref{fig:dice-pmf-cdf}).

정수가 아닌 값에서도 $F$는 정의된다. 계단이 정수 사이에서는 평평하게
이어지므로 $F(3.5) = F(3) = 3/6 = 0.5$다.
\end{example}

\begin{table}[ht]\centering
\begin{tabular}{lrr}
\toprule
$k$ & $P(X=k)$ (확률) & $F(k)=P(X\le k)$ (누적확률) \\
\midrule
1 & 1/6 & 1/6 = 0.167 \\
2 & 1/6 & 2/6 = 0.333 \\
3 & 1/6 & 3/6 = 0.5 \\
4 & 1/6 & 4/6 = 0.667 \\
5 & 1/6 & 5/6 = 0.833 \\
6 & 1/6 & 6/6 = 1 \\
\midrule
$a < 1$ & 없음 & $F(a) = 0$ \\
$a > 6$ & 없음 & $F(a) = 1$ \\
\bottomrule
\end{tabular}
\caption{주사위 눈의 확률(Probability)과 누적확률(Cumulative Probability).}
\label{tab:dice-cdf}
\end{table}

\begin{figure}[ht]\centering
  \includegraphics[width=\linewidth]{04_01_dice_pmf_cdf.png}
  \caption{주사위 눈의 확률질량함수(PMF, 왼쪽)와 누적분포함수(CDF,
  오른쪽). CDF는 각 정수에서 계단으로 뛰어오르고, 뛰어오르기 직전의
  값은 열린 점, 뛰어오른 뒤의 값은 닫힌 점으로 표시된다.}
  \label{fig:dice-pmf-cdf}
\end{figure}

\begin{example}[F(3)과 P(X>3)]
\label{ex:cdf-fifteenths}
확률변수 $X$가 $P(X=x) = x/15$, $x \in \{1,2,3,4,5\}$를 만족한다고
하자. 먼저 이것이 PMF의 조건을 만족하는지 확인한다.

\begin{align*}
  &1+2+3+4+5 = 15               && \text{분자의 합}\\
  &15/15 = 1                    && \text{확률의 합, \Cref{def:pmf}의 조건과 일치}
\end{align*}

$F(3)$은 $x=1,2,3$의 확률을 누적한다.

\begin{align*}
  F(3) &= 1/15 + 2/15 + 3/15 = 6/15 = 0.4
\end{align*}

$P(X>3)$은 여집합으로 구한다.

\begin{align*}
  P(X>3) &= 1 - F(3) = 1 - 0.4 = 0.6
\end{align*}

교차 확인은 $x=4,5$의 확률을 직접 더해서 한다.

\begin{align*}
  P(4) + P(5) &= 4/15 + 5/15 = 9/15 = 0.6
\end{align*}

두 계산이 일치한다.
\end{example}

\begin{pitfall}[이산형에서는 $>$와 $\ge$가 다르다]
\label{pit:cdf-strict}
\Cref{ex:cdf-fifteenths}에서 $P(X>3) = 0.6$이지만 $P(X \ge 3)$은
다르다. $X=3$을 포함하기 때문이다.
\begin{align*}
  P(X \ge 3) &= 1 - F(2) && \text{여집합}\\
  &= 1 - \tfrac{1+2}{15} && \text{$F(2)$를 PMF로 전개}\\
  &= 1 - \tfrac{3}{15} = \tfrac{12}{15} = 0.8 && \text{계산}
\end{align*}
$0.8 \ne 0.6$이다. 이산형에서 경계값을 포함하는지는 결과를 바꾼다.
연속형에서는 $P(X=v)=0$이므로(\Cref{def:pdf}) $>$와 $\ge$의 차이가
사라진다.
\end{pitfall}

\begin{pitfall}[식의 좌우가 다른 변수를 쓰면 안 된다]
\label{pit:cdf-slide-kx}
슬라이드 값은 $P(X>k)=1-P(X\le x)$, 다시 계산하면
$P(X>x)=1-F(x)$, 즉 $P(X>3)=0.6$. 왼쪽 변에 쓴 $k$와 오른쪽 변에
쓴 $x$가 같은 변수를 가리켜야 식이 성립하는데, 슬라이드는 이 둘을
다른 기호로 적었다. 변수 하나를 통일해 $P(X>x) = 1 - F(x)$로 써야
\Cref{def:cdf}와 맞는 식이 된다. 숫자 결과 0.6 자체는 맞다.
\end{pitfall}

\paragraph{중앙값을 CDF로 읽기.} 중앙값의 정의 자리는
\Cref{def:center}에 하나만 둔다. 여기서는 CDF로 중앙값을 찾는 성질만
본다. 중앙값 $m$은
\[
  P(X \le m) \ge \tfrac{1}{2}, \qquad P(X \ge m) \ge \tfrac{1}{2}
\]
을 만족하는 값이다. 주사위에서는 $F(3) = 0.5$이고 $F(4) = 0.667$이므로
$3 \le m \le 4$인 모든 값이 이 조건을 만족하고, 관례로 가운데 값인
3.5를 쓴다. 정규분포(Normal Distribution)처럼 CDF가 매끈하게 증가하는
경우에는 $F(\mu) = 0.5$이므로 중앙값과 평균이 같다.
\Cref{fig:iris-cdfs}의 붓꽃 정규 근사 CDF에서 $(\mu, 0.5) = (5.843,
0.5)$인 점이 바로 이 성질을 보여 준다. 중앙값을 실제로 어떻게
계산하는지는 \Cref{sec:center}에서 이어진다.

붓꽃 데이터에서 서로 다른 두 확률변수의 이론 CDF를 비교할 수 있다
(\Cref{fig:iris-cdfs}). 하나는 꽃받침 길이 자체의 정규
CDF($\mu=5.843$, $\sigma^2=0.681$)이고, 다른 하나는 꽃받침 길이가
아니라 \Cref{ex:iris-long-sepal}의 0/1 판정을 $m=10$번 반복했을 때
``길다''로 판정되는 \textbf{횟수}의 이항(Binomial) CDF($p=0.087$,
$m=10$)다.

\begin{align*}
  F(0) &= 0.913^{10} = 0.4024 && \text{이항, $P(\text{10번 중 0번 ``길다''})$}\\
  F(1) &= 0.7859              && \text{이항, 누적}\\
  F(2) &= 0.9504               && \text{이항, 누적}\\
  F(3) &= 0.9922               && \text{이항, 누적}
\end{align*}

이항분포와 정규분포 각각의 자세한 정의와 모수는
\Cref{sec:binomial}과 \Cref{sec:normal}에서 다룬다.

\begin{figure}[ht]\centering
  \includegraphics[width=\linewidth]{04_03_iris_cdfs.png}
  \caption{붓꽃 데이터에서 나온 이론 CDF 두 개. 정규 CDF($\mu=5.843$,
  $\sigma^2=0.681$, 점 $(\mu, 0.5)$ 표시)는 꽃받침 길이 자체의 분포이고,
  이항 CDF($p=0.087$, $m=10$)는 꽃받침 길이가 아니라 ``길다''로
  판정되는 횟수의 분포다.}
  \label{fig:iris-cdfs}
\end{figure}

\begin{pitfall}[분산은 $n$과 $n-1$ 중 무엇으로 나눴는지가 다르다]
\label{pit:variance-n}
슬라이드 값은 0.681($n$으로 나눔), $n-1$로 다시 계산하면 0.686.
이것은 오류가 아니라 \textbf{관례 차이}다. $n$으로 나누면 모집단
분산의 정의를 그대로 쓴 것이고, $n-1$로 나누면 표본에서 모분산을
추정할 때 생기는 편향(Bias)을 보정한 것이다(\Cref{def:spread},
\Cref{pit:bias}). 이 책의 \Cref{def:spread}는 $n-1$ 쪽 정의를 쓴다.
\end{pitfall}

두 CDF를 겹쳐 놓고 최대 격차를 재는 방법은 \Cref{sec:cdf}에서, 데이터
자체로 계단 모양 CDF를 만드는 경험적 CDF(Empirical CDF)는
\Cref{def:ecdf}에서 다룬다.

% =====================================================================
\section{중심 --- 평균과 중앙값}
\label{sec:center}
% source: deck 4 slide 13, 15-18

\begin{inbrief}
  \item[Why] 데이터 전체를 대표하는 값 하나가 필요하기 때문이다.
  \item[When] \term{평균}(Mean)은 분포가 대칭일 때, \term{중앙값}(Median)은
              한쪽으로 치우쳤을 때 쓴다.
  \item[Where] 소득·가격·반응시간처럼 극단값이 존재하는 데이터에서는 거의 항상
               중앙값 쪽이 안전하다.
  \item[How] 왜도를 먼저 계산해 어느 쪽을 쓸지 정한다(\Cref{sec:skewness}).
\end{inbrief}

\begin{definition}[중심의 두 정의]
\label{def:center}
\term{평균}(Mean)은 모든 값을 더해서 개수로 나눈 값이다.
\[
  \mu = \frac{1}{n}\sum_{i=1}^{n} x_i
\]
\term{중앙값}(Median)은 값을 크기순으로 정렬했을 때 가운데에 오는 값이다.
\end{definition}

두 값의 차이는 계산식이 아니라 \textbf{극단값(Extreme Value)에 대한 반응}에
있다. 평균은 모든 값을 실제로 더하므로, 값 하나가 커지면 그 크기만큼 끌려간다.
중앙값은 ``몇 번째 값인가''만 보므로, 가장 큰 값이 두 배가 되어도 순위가 그대로면
꿈쩍하지 않는다.

이 성질을 \term{저항성}(Resistance)이라고 한다. 중앙값은 저항성이 있고 평균은
없다.

\begin{example}[값 하나가 평균을 움직인다]
\label{ex:resistance}
다섯 개의 값 $\{1,\ 2,\ 3,\ 4,\ 5\}$ 에서 평균과 중앙값은 모두 3이다.
마지막 값만 $5 \to 100$ 으로 바꾸면:

\begin{center}
\begin{tabular}{lrr}
\toprule
데이터 & 평균 & 중앙값 \\
\midrule
$\{1,2,3,4,5\}$   & 3.0  & 3 \\
$\{1,2,3,4,100\}$ & 22.0 & 3 \\
\bottomrule
\end{tabular}
\end{center}

평균은 7배 넘게 뛰었지만 중앙값은 그대로 3이다. 그리고 바뀐 평균 22는
\textbf{다섯 개 값 중 네 개보다 크다} --- 대표값으로서 아무 의미가 없다.
\end{example}

% =====================================================================
\section{퍼짐 --- 표준편차와 IQR}
\label{sec:spread}
% source: deck 4 slide 21-22

\begin{inbrief}
  \item[Why] 중심만으로는 부족하다. 값들이 중심 주변에 모여 있는지 흩어져 있는지를
             알아야 한다.
  \item[When] \term{표준편차}(Standard Deviation)는 대칭 분포에서, \term{사분위
              범위}(Interquartile Range)는 극단값이 있을 때 쓴다.
  \item[Where] 두 값의 차이가 클수록 그 데이터에 극단값이 많다는 신호다.
  \item[How] $\sd$ 와 $\IQR$ 을 \emph{둘 다} 계산해 비교한다. 차이 자체가 정보다.
\end{inbrief}

\begin{definition}[퍼짐의 두 정의]
\label{def:spread}
\term{표준편차}(Standard Deviation)는 각 값이 평균에서 떨어진 거리를 제곱해
평균낸 뒤 제곱근을 취한 값이다.
\[
  \sigma = \sqrt{\frac{1}{n-1}\sum_{i=1}^{n}(x_i - \mu)^2}
\]
\term{사분위 범위}(Interquartile Range)는 가운데 50\%가 차지하는 폭이다.
\[
  \IQR = Q_3 - Q_1
\]
\end{definition}

여기서도 대비는 \Cref{sec:center}와 같다. 거리를 \textbf{제곱}하기 때문에
표준편차는 극단값에 가장 민감하고, 상하위 25\%를 아예 보지 않는 $\IQR$ 은
극단값에 면역이다.

과제 데이터의 \texttt{sugars}(당류) 변수가 이 차이를 극단적으로 보여준다.

\begin{center}
\begin{tabularx}{0.85\linewidth}{@{}l r X@{}}
\toprule
 & 값 & 이 숫자가 실제로 가리키는 것 \\
\midrule
$\sd$  & 21.02 & 밀크셰이크 12개에 끌려 올라간 값 \\
$\IQR$ &  8.00 & 일반 메뉴가 실제로 차지하는 폭 \\
\bottomrule
\end{tabularx}
\end{center}

$\sd = 21.02$ 는 \textbf{실존하지 않는 폭}이다. 대부분의 메뉴는 3--11\,g 사이에
있고, 21이라는 숫자는 소수의 밀크셰이크가 만들어 낸 것이다.

\subsection{IQR이 이상치를 결정한다}
\label{subsec:fence}

박스플롯(Box Plot)에서 점으로 찍히는 \term{이상치}(Outlier)는 눈대중이 아니라
$\IQR$ 로부터 기계적으로 계산된다.

\begin{figure}[ht]\centering
\begin{tikzpicture}[
  node distance = 7mm,
  flowbox/.style = {draw = bookrule, fill = booksurface, rounded corners = 2pt,
                   minimum width = 34mm, minimum height = 9mm, align = center,
                   font = \small},
  arrow/.style  = {-{Stealth[length=2.2mm]}, draw = bookmuted}]
  \node[flowbox] (q)     {$Q_1 = 3$,\quad $Q_3 = 11$};
  \node[flowbox, below = of q]  (iqr)   {$\IQR = 11 - 3 = 8$};
  \node[flowbox, below = of iqr] (fence) {울타리(Fence)\\ $= 11 + 1.5 \times 8 = 23$};
  \node[flowbox, below = of fence, fill = bookwarnbg, draw = bookwarn] (out)
        {23을 넘는 12개\\ = 전부 밀크셰이크};
  \draw[arrow] (q)     -- (iqr);
  \draw[arrow] (iqr)   -- (fence);
  \draw[arrow] (fence) -- (out);
\end{tikzpicture}
\caption{이상치는 $\IQR$ 에서 기계적으로 유도된다. \texttt{sugars} 변수의 실제
계산 과정이다.}
\label{fig:fence-flow}
\end{figure}

\begin{figure}[ht]\centering
  \includegraphics[width=\linewidth]{q1_fig2_iqr_boxplot_anatomy.png}
  \caption{박스플롯의 해부. 상자가 곧 $\IQR$ 이고, 그 $\IQR$ 이
  \Cref{fig:fence-flow}의 울타리를 통해 어떤 점이 이상치인지를 결정한다. x축은
  의도적으로 끊었다 --- 실제 비율로 그리면 상자가 전체 폭의 1/5로 찌그러져 각
  부분에 이름을 붙일 공간이 없다.}
  \label{fig:iqr-anatomy}
\end{figure}

\begin{pitfall}[$1.5 \times \IQR$ 은 법칙이 아니다]
계수 1.5는 통계적으로 유도된 값이 아니라 Tukey가 제안한 \textbf{관례}다
\citep{tukey1977}. 따라서 ``울타리를 넘었으니 이상치다''가 아니라 "이 관례로는
이상치로 분류된다"가 정확한 표현이다. 실제로 \Cref{fig:iqr-anatomy}의 12개는
오류가 아니라 정상적인 밀크셰이크다. \textbf{이상치는 버릴 대상이 아니라 설명할
대상이다.}
\end{pitfall}

% =====================================================================
\section{최빈값과 범위}
\label{sec:mode-range}
% source: deck 4 slide 13, 19-20

\begin{inbrief}
  \item[Why] 평균과 중앙값 말고도, 범주형 데이터에 쓸 수 있는 중심값과
             가장 단순한 퍼짐값이 필요하다.
  \item[When] 값이 범주형이면 최빈값(Mode)을 쓰고, 계산이 빠른 1차
              점검이 필요하면 범위(Range)를 쓴다.
  \item[Where] 최빈값은 명목 척도(Nominal Scale)에도 정의된다
               (\Cref{tab:attr-ops}).
  \item[How] 최빈값은 가장 많이 나온 값을 세고, 범위는 최댓값에서
             최솟값을 뺀다.
\end{inbrief}

\begin{definition}[최빈값과 범위]
\label{def:mode}
확률변수의 \term{최빈값}(Mode)은 PMF 또는 PDF가 가장 큰 값이다. 표본에서는
관측 빈도 $\hat f(x)$가 가장 큰 값을 쓴다.
\[
  \operatorname{mode}(X) = \arg\max_x \hat f(x)
\]
\end{definition}

\begin{definition}[범위]
\label{def:range}
확률변수 $X$의 \term{범위}(Range) $r$은 최댓값과 최솟값의 차이다.
\[
  r = \max\{X\} - \min\{X\}
\]
\end{definition}

\begin{example}[붓꽃 꽃받침 길이의 최빈값과 범위]
\label{ex:mode-range}
꽃받침 길이는 소수 한 자리로 반올림되어 기록되어 있어, 같은 값이
여러 번 나온다. 가장 많이 나온 값은 5.0\,cm로 10송이가 이 값을
가진다(\Cref{fig:iris-mode-range}). 이 값은 \Cref{def:center}의 평균
5.843보다 작다. 범위는

\begin{align*}
  r &= 7.9 - 4.3 = 3.6 \qquad \text{(cm)}
\end{align*}
\end{example}

\begin{figure}[ht]\centering
  \includegraphics[width=\linewidth]{04_04_iris_mode_range.png}
  \caption{붓꽃 꽃받침 길이의 빈도. 최빈값(Mode) 5.0\,cm(10송이)을
  강조하고, 평균 5.843\,cm에 선을 그었으며, 최소 4.3\,cm과 최대
  7.9\,cm 사이의 범위(Range) 3.6\,cm을 화살표로 표시했다.}
  \label{fig:iris-mode-range}
\end{figure}

\Cref{ex:resistance}의 데이터를 다시 쓰면, 최빈값과 범위가 저항성에서
어떻게 갈리는지 보인다. $\{1,2,3,4,5\}$는 모든 값이 한 번씩만 나와
최빈값이 정의되지 않으므로, 값 하나를 중복시킨 $\{1,2,2,3,5\}$로
바꿔서 비교한다.

\begin{center}
\begin{tabular}{lrr}
\toprule
데이터 & 최빈값(Mode) & 범위(Range) \\
\midrule
$\{1,2,2,3,5\}$   & 2 & 4 \\
$\{1,2,2,3,100\}$ & 2 & 99 \\
\bottomrule
\end{tabular}
\end{center}

마지막 값을 $5 \to 100$으로 바꿔도 최빈값은 2로 그대로지만, 범위는
$4 \to 99$로 25배 가까이 뛴다.

\Cref{tab:robustness}은 이 장에서 다룬 여섯 가지 통계량을 저항성
기준으로 정리한다.

\begin{table}[ht]\centering
\begin{tabularx}{\linewidth}{@{}l l X@{}}
\toprule
통계량(Statistic) & 극단값 하나에 흔들리는가 & 근거 \\
\midrule
평균(Mean) & 흔들림 & \Cref{ex:resistance} \\
중앙값(Median) & 안 흔들림 & \Cref{ex:resistance} \\
최빈값(Mode) & 안 흔들림(단, 아래 함정 참고) & \Cref{ex:mode-range} \\
범위(Range) & 흔들림 & \Cref{ex:mode-range} \\
표준편차(Standard Deviation) & 흔들림 & \Cref{sec:spread} \\
사분위범위(IQR) & 안 흔들림 & \Cref{sec:spread} \\
\bottomrule
\end{tabularx}
\caption{여섯 가지 통계량의 저항성(Resistance) 성적표.}
\label{tab:robustness}
\end{table}

\begin{pitfall}[최빈값은 극단값이 아닌 값 하나에도 뒤집힌다]
\label{pit:mode-unstable}
최빈값이 저항성이 있다는 말은 \textbf{극단값}에 대해서만 참이다. 극단값이
\emph{아닌} 값 하나가 바뀌면 최빈값 자체가 뒤집힐 수 있다. $\{1,2,2,3,3,3\}$의
최빈값은 3인데, 극단값이 아닌 값 하나를 $3 \to 2$로 바꾼
$\{1,2,2,2,3,3\}$에서는 최빈값이 2로 바뀐다. 게다가 연속 측정값에서는
모든 값이 한 번씩만 나오면 최빈값이 아예 정의되지 않는다. 붓꽃
꽃받침 길이처럼 소수 한 자리로 반올림되어야 비로소 값이 겹치고
최빈값이 의미를 가진다.
\end{pitfall}

% =====================================================================
\section{비대칭 --- 왜도}
\label{sec:skewness}

\begin{inbrief}
  \item[Why] 평균을 쓸지 중앙값을 쓸지 \emph{객관적으로} 정하기 위해서다.
  \item[When] 요약통계표를 만들 때 항상 함께 계산한다.
  \item[Where] $|\gamma| > 1$ 이면 평균 대신 중앙값을 보고한다.
  \item[How] 부호는 꼬리의 방향, 크기는 치우친 정도를 뜻한다.
\end{inbrief}

\begin{definition}[왜도]
\label{def:skewness}
\term{왜도}(Skewness)는 분포의 비대칭 정도다.
\[
  \gamma = \frac{\E\!\left[(x-\mu)^3\right]}{\sigma^3}
\]
\end{definition}

핵심은 \textbf{세제곱}이다. 제곱이면 부호가 사라지므로 표준편차는 퍼진
\emph{정도}만 잰다. 세제곱은 부호가 살아남는다. 그래서 왜도는 "어느 쪽으로
퍼졌는가"를 잰다.

\begin{pitfall}[왜도(Skewness)는 편향(Bias)이 아니다]
\label{pit:bias}
두 단어가 일상어로는 모두 ``치우침''이지만, 통계학에서는 \textbf{완전히 다른
개념}이고 서로 바꿔 쓸 수 없다.

\begin{center}
\begin{tabularx}{\linewidth}{@{}l l X@{}}
\toprule
용어 & 무엇의 성질인가 & 정의 \\
\midrule
\term{왜도}(Skewness) & \textbf{분포}의 모양 & 분포가 좌우 대칭이 아닌 정도.
  $\gamma = \E[(x-\mu)^3]/\sigma^3$ \\
\term{편향}(Bias) & \textbf{추정량}의 정확도 & 추정값의 기댓값이 참값에서
  벗어난 크기. $\mathrm{Bias}(\hat\theta) = \E[\hat\theta] - \theta$ \\
\bottomrule
\end{tabularx}
\end{center}

왜도는 \emph{데이터가 어떻게 생겼는지}를 말하고, 편향은 \emph{내 추정 방법이
틀렸는지}를 말한다. 완벽하게 편향 없는 방법으로 추정해도 왜도가 큰 분포는
여전히 왜도가 크다. 반대로 완벽하게 대칭인 정규분포에서도 잘못된 추정량을 쓰면
편향이 생긴다.

\textbf{혼동이 실제로 일어나는 지점:} 표준편차의 $n-1$ 나눗셈이 바로 \emph{편향}
보정(Bias Correction)이다. $n$ 으로 나누면 분산이 과소추정되어 편향이 생긴다.
이것은 왜도와 아무 관련이 없다.

번역어로는 \term{왜도} 또는 \term{비대칭도}가 표준이며, \textbf{편향은 오역}이다.
\end{pitfall}

\subsection{실제로 계산하는 법 --- $g_1$ 과 $G_1$}
\label{sec:skew-compute}

\Cref{def:skewness}의 $\gamma$ 는 모집단(Population)의 정의다. 손에 있는 것은 표본이므로
실제로 계산하는 값은 두 단계를 거친다.

\textbf{1단계. 보정 없는 표본 왜도 $g_1$.} 정의식의 기댓값을 표본 평균으로 바꾼다.
$k$차 \term{중심적률}(Central Moment)을 $m_k = \frac{1}{n}\sum_{i=1}^{n}(x_i - \bar{x})^k$
라 하면
\[
  g_1 = \frac{m_3}{m_2^{3/2}} .
\]

\textbf{2단계. 표본 보정 $G_1$.} $g_1$ 은 표본에서 왜도를 과소추정하므로 계수를 곱한다.
\[
  G_1 = \frac{\sqrt{n(n-1)}}{n-2}\; g_1 .
\]
이 $G_1$ 이 \term{조정된 Fisher-Pearson 표준화 적률계수}(Adjusted Fisher-Pearson
Standardized Moment Coefficient)이고, \texttt{pandas} 의 \texttt{.skew()} 가 돌려주는 값이다.

손으로 구하는 절차는 다섯 줄이다.

\begin{enumerate}
  \item 평균 $\bar{x}$ 를 구한다.
  \item 편차의 제곱을 모두 더해 $n$ 으로 나눈다. 이것이 $m_2$ 다.
  \item 편차의 세제곱을 모두 더해 $n$ 으로 나눈다. 이것이 $m_3$ 다.
  \item $g_1 = m_3 / m_2^{1.5}$ 를 계산한다.
  \item $\sqrt{n(n-1)}/(n-2)$ 를 곱해 $G_1$ 을 얻는다.
\end{enumerate}

\begin{example}[다섯 개 값으로 끝까지 계산하기]
\label{ex:skew-by-hand}
데이터 $\{1,\ 2,\ 3,\ 4,\ 10\}$, $n = 5$ 이다.

\begin{center}
\begin{tabular}{lrrr}
\toprule
$x_i$ & $x_i - \bar{x}$ & $(x_i-\bar{x})^2$ & $(x_i-\bar{x})^3$ \\
\midrule
1  & $-3$ & 9  & $-27$ \\
2  & $-2$ & 4  & $-8$  \\
3  & $-1$ & 1  & $-1$  \\
4  & $0$  & 0  & $0$   \\
10 & $6$  & 36 & $216$ \\
\midrule
합 & $0$ & 50 & 180 \\
\bottomrule
\end{tabular}
\end{center}

평균은 $\bar{x} = 20/5 = 4$ 다. 이어서

\begin{align*}
  m_2 &= 50/5 = 10, \qquad m_3 = 180/5 = 36, \\
  g_1 &= \frac{36}{10^{1.5}} = \frac{36}{31.6228} = 1.1384, \\
  G_1 &= \frac{\sqrt{5 \cdot 4}}{3}\, g_1 = 1.4907 \times 1.1384 = 1.6971 .
\end{align*}

\texttt{pd.Series([1,2,3,4,10]).skew()} 가 정확히 $1.6971$ 을 돌려준다.

편차 세제곱 열을 보면 왜도가 어디서 오는지 한눈에 보인다. 값 10 하나가 $+216$ 을
만들고 나머지 넷이 합쳐서 $-36$ 밖에 못 만든다. \textbf{세제곱이 먼 값에 압도적인
가중치를 준다.}
\end{example}

보정 계수 $\sqrt{n(n-1)}/(n-2)$ 는 $n$ 이 커지면 1로 수렴한다. \Cref{ex:skew-by-hand}
처럼 $n=5$ 면 1.49로 크지만, 이 책의 데이터처럼 $n=126$ 이면 1.0121이라 $g_1$ 과 $G_1$ 의
차이는 1\% 남짓이다.

\begin{center}
\begin{tabular}{lr}
\toprule
\texttt{sugars} 의 왜도 계산 방식 & 값 \\
\midrule
$g_1$ (보정 전)                                  & 2.6765 \\
$G_1$ (보정 후)                                  & 2.7088 \\
$m_3 / s^3$ (분모만 $\mathrm{ddof}=1$)           & 2.6447 \\
\texttt{pandas .skew()}                          & 2.7088 \\
\bottomrule
\end{tabular}
\end{center}

\begin{pitfall}[분자와 분모의 나눗셈을 섞지 않는다]
위 표의 셋째 줄 $m_3/s^3$ 는 흔한 오답이다. 분모에는 표본표준편차($\mathrm{ddof}=1$,
$n-1$ 로 나눔)를 쓰면서 분자 $m_3$ 는 $n$ 으로 나눈 값이라, \textbf{어느 정의와도
일치하지 않는다}. $m_2$ 와 $m_3$ 는 둘 다 $n$ 으로 나누고, 보정은 마지막에 계수로
한 번만 한다.
\end{pitfall}

\subsection{부호와 방향}
\label{sec:skew-sign}

\begin{center}
\begin{tabularx}{\linewidth}{@{}l l l X@{}}
\toprule
부호 & 이름 & 모양 & 평균과 중앙값의 관계 \\
\midrule
$\gamma > 0$ & 우측 왜도(Right-skewed) & 오른쪽 꼬리가 김 & 평균 $>$ 중앙값 \\
$\gamma = 0$ & 대칭(Symmetric)         & 종 모양         & 평균 $=$ 중앙값 \\
$\gamma < 0$ & 좌측 왜도(Left-skewed)  & 왼쪽 꼬리가 김  & 평균 $<$ 중앙값 \\
\bottomrule
\end{tabularx}
\end{center}

왜 평균만 끌려가는지는 \Cref{ex:resistance}에서 이미 봤다. 그래서 왜도는
그림에서 \textbf{두 세로선 사이의 틈}으로 나타난다.

\begin{figure}[ht]\centering
  \includegraphics[width=\linewidth]{q1_fig1_skewness_mean_vs_median.png}
  \caption{왜도를 눈으로 보는 법. 왜도가 커질수록($0.24 \to 2.71$) 중앙값(검정
  실선)과 평균(주황 점선) 사이의 틈이 벌어진다.}
  \label{fig:skew-gap}
\end{figure}

\begin{pitfall}[틈은 반드시 SD 단위로 비교한다]
원래 단위로 재면 거짓말이 된다. \texttt{sodium}(나트륨)의 틈은 43.75\,mg이고
\texttt{sugars}(당류)의 틈은 6.29\,g이지만, 왜도는 0.42 대 2.71로 정반대다.
단위가 다르기 때문이다. $(\mu - \mathrm{median})/\sigma$ 로 표준화해야
0.08 대 0.30이 되어 왜도와 순서가 맞는다. \Cref{fig:skew-gap}의 각 패널에
적힌 \texttt{gap} 값이 이 표준화된 값이다.
\end{pitfall}

같은 표본을 강의 5강(Basic Statistics and Visualization)에서 다룬 다섯 가지 그림,
즉 히스토그램(Histogram, \Cref{sec:histogram}), 선 그래프(Line Chart, 밀도 곡선
Density Line, \Cref{sec:density-line}), 박스플롯(Box Plot, \Cref{sec:boxplot}),
바이올린 플롯(Violin Plot, \Cref{sec:violin}), 바코드 차트(Barcode Chart,
\Cref{sec:barcode})로 나란히 그리면, 왜도의 부호가 그림마다 어떤 흔적을 남기는지
한눈에 비교할 수 있다(\Cref{tab:skew-by-chart}). \Cref{fig:skew-right}부터
\Cref{fig:skew-left}까지 세 그림의 데이터는
\textbf{실제 강의 자료가 아니라 합성(synthetic) 표본}이다. 우측 왜도
표본은 고정 시드(seed)로 뽑은 감마(Gamma) 분포 난수이고, 좌측 왜도 표본은 그 값을
$60$에서 뺀 거울상($60-x$)이라 $G_1$ 의 부호가 정확히 뒤집히며, 대칭 표본은 평균과
표준편차를 맞춘 정규(Normal) 분포 난수다. 각 패널의 검정 실선은 중앙값(median),
주황 점선은 평균(mean)으로 \Cref{fig:skew-gap}과 같은 인코딩이다.

\begin{example}[세 표본에서 왜도의 부호를 손으로 확인하기]
\label{ex:skew-three-cases}
세 표본 모두 $n=300$ 이다. 순서는 (a) $G_1$ 의 부호를 본다, (b) 그 부호가
평균$-$중앙값의 부호로 이어지는지 확인한다, (c) 그 부호로부터 긴 꼬리가 어느
쪽인지 결정한다.

\textbf{우측 왜도 표본.} $G_1 = 1.48 > 0$ 이므로 평균이 중앙값보다 커야 한다.
\begin{align*}
  &16.32 - 13.95 = 2.37       && \text{평균 $-$ 중앙값 (양수, 예상과 일치)}\\
  &2.37 / 11.15 = 0.21        && \text{SD 단위로 표준화한 틈}\\
  &22.06 - 13.95 = 8.11       && \text{$Q_3 - $ 중앙값}\\
  &13.95 - 8.40 = 5.55        && \text{중앙값 $- Q_1$}
\end{align*}
위쪽 거리(8.11)가 아래쪽 거리(5.55)보다 커서, 중앙값 위로 값이 더 퍼져 있다.
즉 긴 꼬리는 \textbf{오른쪽}이다. $G_1$ 의 부호, 평균$-$중앙값의 부호, 꼬리의
방향이 셋 다 같은 이야기를 하고 있다.

\textbf{좌측 왜도 표본.} $G_1 = -1.48 < 0$ 이므로 평균이 중앙값보다 작아야 한다.
\begin{align*}
  &43.68 - 46.05 = -2.37      && \text{평균 $-$ 중앙값 (음수, 예상과 일치)}\\
  &-2.37 / 11.15 = -0.21      && \text{SD 단위로 표준화한 틈}\\
  &51.60 - 46.05 = 5.55       && \text{$Q_3 - $ 중앙값}\\
  &46.05 - 37.94 = 8.11       && \text{중앙값 $- Q_1$}
\end{align*}
이번에는 아래쪽 거리(8.11)가 위쪽 거리(5.55)보다 커서, 긴 꼬리는
\textbf{왼쪽}이다. 좌측 왜도 표본은 우측 왜도 표본을 $60-x$ 로 뒤집어 만들었으므로,
틈($0.21 \leftrightarrow -0.21$)과 두 분위수 거리($8.11 \leftrightarrow 5.55$)가
정확히 자리를 바꾼다. 이것이 \Cref{fig:skew-left}의 표본이 \Cref{fig:skew-right}의 표본을
좌우로 뒤집은 모습으로 보이는 이유다.

\textbf{대칭 표본.} $G_1 = -0.07$ 로 부호는 음수지만 크기가 매우 작다.
\begin{align*}
  &16.48 - 16.78 = -0.30      && \text{평균 $-$ 중앙값}\\
  &-0.30 / 10.95 = -0.03      && \text{SD 단위로 표준화한 틈}\\
  &23.49 - 16.78 = 6.71       && \text{$Q_3 - $ 중앙값}\\
  &16.78 - 9.55 = 7.23        && \text{중앙값 $- Q_1$}
\end{align*}
두 분위수 거리(6.71과 7.23)가 서로 가까워서 어느 쪽으로도 뚜렷한 꼬리가
없다. $G_1$ 이 정확히 $0$ 이 아닌 이유는 유한한 표본이 완벽하게 대칭일 수는
없기 때문이다. $|G_1|=0.07 < 0.5$ 이므로 \Cref{sec:skew-rule}의 기준으로도
``거의 대칭''으로 분류된다.
\end{example}

\begin{figure}[p]\centering
  \includegraphics[width=\linewidth,height=0.85\textheight,keepaspectratio]{02_skew_right.png}
  \caption{우측 왜도(Right-skewed, $G_1=1.48$) 표본을 다섯 가지 그림으로 본 모습.
  히스토그램(Histogram)과 선 그래프(Line Chart, 밀도 곡선 Density Line)는 왼쪽에
  봉우리가 몰리고 오른쪽으로 긴 꼬리가 남으며, 박스플롯(Box Plot)은
  중앙값(Median, 검정 실선)이 상자 왼쪽으로 치우치고 오른쪽에
  이상치(Outlier, 주황 점)가 늘어선다. 바이올린 플롯(Violin Plot)은 오른쪽으로
  갈수록 가늘어지고, 바코드 차트(Barcode Chart)는 막대가 왼쪽에 몰렸다가
  오른쪽으로 갈수록 성겨진다.}
  \label{fig:skew-right}
\end{figure}

\begin{figure}[p]\centering
  \includegraphics[width=\linewidth,height=0.85\textheight,keepaspectratio]{02_skew_symmetric.png}
  \caption{대칭(Symmetric, $G_1=-0.07$) 표본을 다섯 가지 그림으로 본 모습. 중앙값과
  평균이 거의 겹치고, 박스플롯의 두 수염과 바이올린 플롯의 양쪽 폭, 바코드 차트의
  양쪽 밀도가 서로 비슷하다.}
  \label{fig:skew-symmetric}
\end{figure}

\begin{figure}[p]\centering
  \includegraphics[width=\linewidth,height=0.85\textheight,keepaspectratio]{02_skew_left.png}
  \caption{좌측 왜도(Left-skewed, $G_1=-1.48$) 표본을 다섯 가지 그림으로 본 모습.
  \Cref{fig:skew-right}를 좌우로 뒤집은 모양과 같다. 히스토그램(Histogram)과
  선 그래프(Line Chart, 밀도 곡선 Density Line)는 오른쪽에 봉우리가 몰리고
  왼쪽으로 긴 꼬리가 남으며, 박스플롯(Box Plot)은 중앙값(Median)이 상자
  오른쪽으로 치우치고 왼쪽에 이상치(Outlier)가 늘어선다.}
  \label{fig:skew-left}
\end{figure}

\begin{table}[htbp]\centering
\begin{tabularx}{\linewidth}{@{}l X X X@{}}
\toprule
그림(Chart) & 우측 왜도(Right-skewed) & 대칭(Symmetric) & 좌측 왜도(Left-skewed) \\
\midrule
히스토그램(Histogram) & 봉우리가 왼쪽, 막대가 오른쪽으로 길게 늘어짐
  & 봉우리가 가운데, 양쪽 막대 길이가 비슷함
  & 봉우리가 오른쪽, 막대가 왼쪽으로 길게 늘어짐 \\
선 그래프(Line Chart, 밀도 곡선 Density Line) & 봉우리가 평균 왼쪽, 오른쪽 꼬리가 김
  & 봉우리가 평균과 거의 겹침
  & 봉우리가 평균 오른쪽, 왼쪽 꼬리가 김 \\
박스플롯(Box Plot) & 중앙값 선이 상자 왼쪽에 치우침, 오른쪽 수염이 길고
  오른쪽에 이상치
  & 중앙값 선이 상자 가운데 근처, 양쪽 수염 길이가 비슷함
  & 중앙값 선이 상자 오른쪽에 치우침, 왼쪽 수염이 길고 왼쪽에 이상치 \\
바이올린 플롯(Violin Plot) & 오른쪽으로 가늘고 긴 꼬리
  & 좌우가 거의 대칭인 모양
  & 왼쪽으로 가늘고 긴 꼬리 \\
바코드 차트(Barcode Chart) & 막대가 왼쪽에 몰리고 오른쪽으로 갈수록 성겨짐
  & 막대가 가운데를 중심으로 고르게 퍼짐
  & 막대가 오른쪽에 몰리고 왼쪽으로 갈수록 성겨짐 \\
\bottomrule
\end{tabularx}
\caption{다섯 가지 그림(Chart)에서 왜도(Skewness)가 남기는 흔적.}
\label{tab:skew-by-chart}
\end{table}

\begin{pitfall}[박스플롯에는 평균이 기본으로 그려지지 않는다]
박스플롯은 원래 $Q_1$과 $Q_3$ 사이의 상자, 중앙값 선, 그리고 $1.5\times\IQR$
이내의 가장 먼 관측값까지 뻗는 수염만 그리며(\Cref{sec:boxplot}), 그 밖의 값은
점으로 찍히는 이상치다. 평균은 기본값으로 표시되지 않는다. \Cref{fig:skew-right}와
\Cref{fig:skew-left}에서는 설명을 위해 평균 점선을 일부러 덧그린 것이다. 평균 선이
없는 보통의 박스플롯에서 왜도를 읽으려면, 중앙값 선이 상자의 가운데에서 어느 쪽으로
치우쳤는지, 두 수염의 길이 차이, 그리고 이상치가 어느 쪽에 몰려 있는지를 본다.
\end{pitfall}

\subsection{판단 기준}
\label{sec:skew-rule}

\begin{center}
\begin{tabular}{@{}ll@{}}
\toprule
$|\gamma| < 0.5$        & 거의 대칭 --- 평균 + 표준편차 \\
$0.5 \le |\gamma| \le 1$ & 중간 --- 평균도 쓸 만하나 주의 \\
$|\gamma| > 1$          & 강한 왜도 --- \textbf{중앙값 + $\IQR$} \\
\bottomrule
\end{tabular}
\end{center}

\texttt{sugars} 는 $\gamma = 2.71$ 이고, 결정적으로 \textbf{평균 13.29가
$Q_3 = 11$ 보다 크다}. 즉 전체의 82.5\%가 ``평균 이하''다. 평균이라는 단어의
일상적 의미가 깨진 상태다(\Cref{ex:resistance}의 22와 같은 상황이다).

\begin{pitfall}[왜도의 정의는 하나가 아니다]
과목에 따라 \term{Pearson 제2 왜도계수}(Pearson's Second Skewness Coefficient)를
가르치기도 한다.
\[
  \mathrm{Sk}_2 = \frac{3(\bar{x} - \mathrm{median})}{s}
\]
같은 데이터에서 이 값은 \texttt{serving\_size} 0.19, \texttt{calories} 0.21,
\texttt{sodium} 0.25, \texttt{sugars} 0.90 이다. 부호는 $G_1$ 과 같지만 크기가 전혀
다르다. \texttt{sugars} 가 2.71이 아니라 0.90으로 나오므로, 위의 ``$|\gamma| > 1$ 이면
중앙값'' 기준을 그대로 적용하면 \textbf{판정이 뒤집힌다}. 두 지표는 스케일이 다르니
임계값을 섞어 쓰면 안 된다.
\end{pitfall}

\begin{aside}
R로 교차검증한다면 패키지마다 정의가 다르다는 점을 예상해야 한다.
\texttt{e1071::skewness(type = 2)} 가 \texttt{pandas} 와 같은 $G_1$ 이고,
\texttt{moments::skewness()} 는 보정 없는 $g_1$ 을 돌려준다. \texttt{sugars} 에서
2.7088과 2.6765의 차이로 나타난다.
\end{aside}

\begin{pitfall}[왜도는 모양을 알려 주지 않는다]
$\gamma = 2.71$ 만 보면 ``완만하게 늘어진 긴 꼬리''를 상상하게 된다. 실제로는
0--20\,g에 114개, 56--93\,g에 12개가 있고 \textbf{그 사이 20--56\,g은 완전히
비어 있다}. 왜도는 ``평균을 쓸까 중앙값을 쓸까''를 정해 주는 스위치일 뿐,
분포의 모양은 반드시 그림으로 확인해야 한다 --- \analogy{지도}로는 보이지 않고
\analogy{항공사진}이 필요한 지점이다.
\end{pitfall}

\begin{figure}[ht]\centering
  \includegraphics[width=\linewidth]{q1_fig3_skew_cause_two_populations.png}
  \caption{왜도 2.71의 정체. 긴 꼬리가 아니라 구멍을 사이에 둔 두 집단
  (Mixture of Two Populations)이다.}
  \label{fig:two-populations}
\end{figure}

% =====================================================================
\section{결측값 --- 0은 결측이 아니다}
\label{sec:missing}

\begin{inbrief}
  \item[Why] 0과 결측을 섞으면 계산은 조용히 성공하고 결과만 틀린다.
  \item[When] 데이터를 처음 받아 \texttt{isna()} 를 돌린 직후에 판단한다.
  \item[Where] 센서 데이터, 설문 응답, 의료 기록처럼 0이 ``없음''의 코드로 쓰이는
               분야에서 특히 위험하다.
  \item[How] ``이 변수에서 0이 \emph{물리적으로 가능한 값인가}''를 묻는다.
\end{inbrief}

\term{결측값}(Missing Value)은 \textbf{값이 기록되지 않았다}는 뜻이다. 0은
\textbf{측정했더니 0이었다}는 뜻이다. 둘은 완전히 다른 정보이고, 섞으면 복구할
수 없다.

\begin{example}[0을 결측으로 처리하면]
\label{ex:zero-as-missing}
과제 데이터에서 \texttt{sugars == 0} 인 항목은 17개이고, 전부 감자튀김(10개)과
치킨너겟(7개)이다. 이들은 실제로 당류가 없는 음식이다. 이를 결측으로 바꾸면:

\begin{center}
\begin{tabular}{@{}lrrrr@{}}
\toprule
 & $n$ & 평균 & 중앙값 & $\IQR$ \\
\midrule
0을 값으로 (옳음)   & 126 & 13.29 & 7.0 & 8.0 \\
0을 결측으로 (틀림) & 109 & 15.36 & 8.0 & 6.0 \\
\bottomrule
\end{tabular}
\end{center}

감자튀김의 당류 중앙값이 0\,g에서 1.15\,g으로 바뀐다. \textbf{당류가 없는
음식이 당류를 가진 것처럼 보이게 된다.} 게다가 $n$ 이 126에서 109로 줄어
``결측 없는 데이터''라는 사실 자체가 사라진다.
\end{example}

\begin{keyidea}
판단 기준은 하나다. \textbf{그 0이 물리적으로 가능한 값인가?}

\begin{itemize}
  \item 당류 0\,g, 강수량 0\,mm, 사고 건수 0건 $\rightarrow$ \textbf{가능하다. 값이다.}
  \item 혈압 0, 키 0\,cm, 성인 나이 0세 $\rightarrow$ \textbf{불가능하다.} 결측을
        0으로 채워 넣은 흔적일 가능성이 높다.
\end{itemize}

후자의 경우에만 \verb|df.replace(0, np.nan)| 이 정당하다.
\end{keyidea}

% =====================================================================
\section{어떻게 적용하는가}
\label{sec:desc-apply}

이 장의 도구는 과제와 다른 장에서 바로 쓰인다.

\begin{itemize}
  \item HW1의 버거 대 밀크셰이크 비교(\Cref{tab:hw1-burger-vs-shake})는
        이 장의 평균과 퍼짐(표준편차, 최솟값, 최댓값)으로 두 음식 유형을 비교한 결과다.
        중앙값과 $\IQR$은 HW1 표에 나오지 않는다.
  \item HW2의 추정량(\Cref{sec:hw2-estimators})은
        \Cref{subsec:expectation}의 기댓값을 표본에서 추정하는
        방법을 다룬다. HW2의 분포 비교(\Cref{sec:hw2-compare-how})는
        이 장의 CDF를 겹쳐 보는 방식을 쓴다.
  \item 경험적 CDF로 후보 분포를 눈으로 확인하는 절차는
        \Cref{sec:dist-check-cdf}에서 이어진다.
  \item \Cref{ex:iris-long-sepal}의 0/1 확률변수 $A$와 같은 발상은
        이산화(Discretization, \Cref{sec:discretization})로 일반화된다.
  \item 결측값을 최빈값으로 채우는 방법은 \Cref{sec:impute}에서 다룬다.
  \item 히스토그램(\Cref{sec:histogram})은 표본의 경험적 PDF고,
        분위수(\Cref{sec:quantile})는 CDF의 역함수다.
\end{itemize}

파이썬 \texttt{pandas}로는 이 절의 값들을 각각 \texttt{.mode()},
\texttt{.max() - .min()}, \texttt{(x >= 7).mean()}으로 구한다. 마지막
식의 결과는 $13/150 = 0.0867$로 \Cref{sec:pmf-pdf}의 $f(1)$과 같다.

% =====================================================================
\section{앞으로}
\label{sec:desc-future}

이 장에서 다룬 확률변수, PMF/PDF, CDF는 이름 있는 분포를 다루는 다음
장으로 이어진다(\Cref{ch:distributions}). 특히 이항분포와 정규분포는
\Cref{sec:binomial}과 \Cref{sec:normal}에서 모수를 추정하고 적합도를
검정하는 법까지 다룬다.

두 CDF를 겹쳐 최대 격차를 재는 방법(\Cref{sec:cdf})은 콜모고로프--스미르노프
검정(Kolmogorov--Smirnov Test, \Cref{sec:cd-ks})으로 정식화된다.

표본 평균이 모평균을 얼마나 잘 맞히는지, 즉 \term{불편추정량}(Unbiased
Estimator)이라는 성질이 실제로 무엇을 보장하는지는 가설검정
(\Cref{ch:hypothesis})에서 이어진다.

% =====================================================================
\section{요약}

\begin{itemize}
  \item 평균과 표준편차는 극단값에 끌려가고, 중앙값과 $\IQR$ 은 저항한다
        (\Cref{def:center}, \Cref{def:spread}).
  \item 왜도가 어느 쪽을 쓸지 결정한다. $|\gamma| > 1$ 이면 중앙값이다.
  \item 왜도(Skewness)와 편향(Bias)은 다른 개념이다(\Cref{pit:bias}).
  \item 왜도는 ``얼마나 치우쳤나''만 알려 준다. \emph{왜} 치우쳤는지는 그림의 몫이다
        (\Cref{fig:two-populations}).
  \item 0과 결측은 다르다. 0이 물리적으로 가능한지를 먼저 묻는다
        (\Cref{ex:zero-as-missing}).
\end{itemize}
